\begin{papercard}{mehrnezhad2024}
\cf{Stages} \stagetag{H} (FemTech threat actors and data flows for the fertility harm tree), \stagetag{S} (how EU, UK and Swiss law treat FemTech data), \stagetag{D} (tracker and permission checks of apps and websites). Frontiers in the Internet of Things, 2024 (Newcastle, Royal Holloway, ETH Zurich); Mehrnezhad is a leading FemTech security and privacy researcher. Added from the second search round of 23 October 2026; the CHI 2021 paper by Mehrnezhad and Almeida is still to be obtained.

\cf{Aim and method} Review the GDPR, the Swiss data protection act, the UK MHRA rules and the EU Medical Devices Regulation for FemTech, and test devices and apps for exploitative practices \pc{mehrnezhad2024}{1}{We run experiments on a range of FemTech devices and apps and identify several exploitative practices}. The laws' public documents were searched for FemTech data categories by keyword \pc{mehrnezhad2024}{4}{For each law, we go through its public documents and manually search for mentions of Fem-Tech data via a few keywords}. Websites were opened in Brave to count trackers before any engagement with the cookie notice \pc{mehrnezhad2024}{5}{to identify how many trackers are activated when the website is loaded for the first time, and before any engagement with the cookie notice}; apps were checked with Exodus for trackers and permissions.

\cf{Findings} FemTech data are not named in the GDPR; they overlap with special category data \pc{mehrnezhad2024}{6}{When we search in the GDPR articles and guidelines, Fem-Tech data categories are not mentioned directly}. The UK medical-device regulation is over 20 years old and silent on security and privacy \pc{mehrnezhad2024}{7}{This regulation is more than 20 years old and does not provide any content on the SP aspects of modern medical devices}. Apps had 1 to 9 trackers, and most websites tracked users before they responded to the cookie notice \pc{mehrnezhad2024}{10}{the majority of these websites are tracking the users before the user engages with the cookie notice}. Data concern not only the user but partners, children and others \pc{mehrnezhad2024}{2}{collecting a wide range of sensitive data about the users and others such as partner(s), baby/child, family and friends}.

\cf{Threat actors} From earlier work by the group: \pc{mehrnezhad2024}{3}{These interested parties include, but are not limited to: (ex-)partner and family, employers and colleagues, insurance firms, advertising companies, political and religious organisations, governments, and medical and research companies}.

\cf{Limits (card)} Products chosen to illustrate, not sampled; few apps; UK, EU and Swiss law as of 2023; tracker counts depend on the tools used.

\cf{Use for this thesis} (1) Fertility harm tree: a published list of adversaries (including intimate partners and employers) that goes beyond the post-Roe state actor and fits a German setting. (2) Harms to third parties (partner, child) are a coverage question for Axis A and for Citron and Solove's relationship harms \pp{citron2021}. (3) Evidence that EU law does not name FemTech data, which frames how German providers can present such data as ordinary ``health and fitness'' data. (4) The Brave and Exodus checks are a cheap way to add an observed-practice note to the walkthrough without network analysis.
\end{papercard}
